% ============================================================
% MCT-L1 SECTION — for inclusion in Visioni et al. GeoMIP7
% Tier 2 > Process understanding experiments
% Draft: 30 March 2026 | D. Villanueva et al.
% ============================================================
%
% TABLE ROW (replace placeholder in tab:placeholder):
%   MCT-L1 & INP seeding poleward of 60$^\circ$ (3 levels + negative batch) & Fixed SST & min 5\,yr \\
%

\subsubsection{MCT-L1: Mixed-Phase Cloud Thinning experiment}
\label{sec:mct}

Mixed-Phase Cloud Thinning (MCT) is a proposed radiation
modification method that exploits the thermodynamic instability of
mixed-phase clouds — clouds containing supercooled liquid droplets coexisting
with ice crystals at temperatures between approximately 0 and $-39\,^\circ$C.
By seeding these clouds with ice-nucleating particles (INPs), liquid droplets
are converted to ice via heterogeneous freezing, triggering the
Wegener-Bergeron-Findeisen (WBF) process that further depletes liquid water
\citep{Villanueva2025Science}.
The resulting glaciated cloud has a lower optical depth and shorter lifetime,
reducing its longwave cloud radiative effect; over the polar oceans in winter,
where mixed-phase clouds exert a net longwave warming that dominates over
shortwave cooling, MCT thus produces a net surface cooling
\citep{Villanueva2022ERL}.
Despite growing interest in MCT as a polar climate intervention, no
coordinated multi-model protocol had previously been proposed; here we report
preliminary results from a first multi-model intercomparison effort —
MCT-MIP Level~0 — and outline the CMIP7-aligned Level~1 protocol (MCT-L1).

\paragraph{Level 0: Exploratory multi-model results.}

The Level~0 protocol was designed to be broadly implementable: atmosphere-only
simulations with prescribed, fixed sea-surface temperatures and sea-ice
fractions, perturbing the polar ocean domain (including sea-ice) poleward of 60$^\circ$N and
60$^\circ$S with spatially uniform, time-constant additions to the background
dust concentration at three INP levels ($10^4$, $10^5$, and $10^6$\,L$^{-1}$),
alongside an unperturbed control.
Simulations used present-day (year 2000; repeating if possible) boundary conditions and ran for a
minimum of 5 years (2 years for nudged models), with a 1-year spin-up
discarded for free-running integrations.
Four models participated (Table~\ref{tab:mct_models}): ECHAM-HAM, ICON-HAM,
CESM2, and E3SM.

\begin{table}[h]
\centering
\footnotesize
\begin{tabular}{lll}
\hline
                   & \textbf{ECHAM-HAM}               & \textbf{ICON-HAM}               \\
\hline
Resolution         & T63 ($\sim$1.875$^\circ$), 96$\times$192  & R2B4 ($\sim$160\,km), 20{,}480 cells  \\
Vertical levels    & 31                                & 47                              \\
Grid type          & Gaussian                          & Icosahedral                     \\
Microphysics       & P3, N12                           & 2-mom., N12 \& L06              \\
Run length         & 25\,yr (MLO) \& 2\,yr (nudged)    & 5\,yr (nudged)                  \\
\hline
                   & \textbf{CESM2}                    & \textbf{E3SM v3}                \\
\hline
Resolution         & 0.9$^\circ$, 192$\times$288       & Ne30 ($\sim$1$^\circ$), 21{,}600 cols \\
Vertical levels    & 32                                & 80                              \\
Grid type          & Regular lat-lon                   & Cubed-sphere                    \\
Microphysics       & 2-mom., Wang et al.\ 2014         & P3, Wang et al.\ 2014           \\
Run length         & 5\,yr                             & 5\,yr                           \\
\hline
\end{tabular}
\caption{Models participating in MCT-MIP Level~0, with their atmospheric grid and run
configuration. MLO: Mixed-Layer Ocean. N12: \citet{Niemand2012}; L06: \citet{Lohmann2006}.}
\label{tab:mct_models}
\end{table}

The spatial distribution of the cloud radiative effect (CRE) response during
the November--February season (NDJF) is shown in
Figure~\ref{fig:dcre_maps} for all four models at the mid INP level
($10^2$\,cm$^{-3}$).
Three of the four models produce a net cooling with increasing INP
concentration, consistent with the expected MCT mechanism; E3SM is an outlier,
showing a warming response.

\begin{figure}[h]
\centering
\includegraphics[width=0.85\textwidth]{../figures/dcre_maps_NDJF.png}
\caption{Spatial maps of the net cloud radiative effect perturbation
($\Delta$CRE, W\,m$^{-2}$) during NDJF for the mid INP level
($10^2$\,cm$^{-3}$) relative to the control simulation, for all four
MCT-MIP Level~0 models.
Numbers inset in each panel give the area-weighted mean $\Delta$CRE over
30--90$^\circ$N, 60--90$^\circ$N, and the outlined target region
(60--90$^\circ$N, 30$^\circ$W--90$^\circ$E, i.e.\ the North Atlantic--European sector).
Blue shading indicates cooling; red shading indicates warming.
ECHAM-HAM, CESM2, and ICON-HAM all produce a net negative $\Delta$CRE
poleward of 60$^\circ$N, consistent with the MCT mechanism, while E3SM
shows a widespread warming response driven by a base state dominated by
very low liquid water content.}
\label{fig:dcre_maps}
\end{figure}

The Arctic-mean (60--90$^\circ$N) NDJF responses of CRE, ice water path
(IWP), liquid water path (LWP), and total water path (TWP; from the
Cloud\_cci v3 product, \citealt{Stengel2020}) as a function of INP
concentration are shown in Figure~\ref{fig:arctic_overlay}.
CESM2 and ICON-HAM show a clear reduction in LWP with increasing INP
concentration, consistent with the MCT mechanism, while E3SM maintains very
low LWP across all perturbation levels.
Note that the ECHAM-HAM (MLO) configuration is shown for reference only;
this idealized perturbation \citep{Villanueva2022ERL} was designed to maximize
the droplet freezing rate and is not directly comparable to the
dust-based INP seeding protocol used by the other simulations.

Investigation of the E3SM case reveals a base state in which cloud ice
strongly dominates over liquid (Figure~\ref{fig:arctic_overlay}b--c),
suggesting that clouds in the perturbed domain are already in a
saturated-ice regime where the WBF process cannot operate effectively and
additional INPs likely suppress the WBF process even further, by decreasing the
size of ice particles and slowing down riming and sedimentation.
This result underlines the importance of constraining the mixed-phase cloud
liquid/ice partition against observations in future intercomparisons.

The sensitivity of the cooling to INP concentration also varies substantially
across the three responding models, and the lowest perturbation level
($10^4$\,L$^{-1}$) already produces a strong response in total water path
(Figure~\ref{fig:arctic_overlay}c),
motivating the extension to lower concentrations in the next Level~1 intercomparison.

\begin{figure}[h]
\centering
\captionsetup[subfigure]{position=top, singlelinecheck=false}
% --- top row ---
\begin{subfigure}[b]{0.48\textwidth}
  \caption{}
  \includegraphics[width=\textwidth]{../figures/cre_arctic_NDJF_overlay.png}
\end{subfigure}
\hfill
\begin{subfigure}[b]{0.48\textwidth}
  \caption{}
  \includegraphics[width=\textwidth]{../figures/iwp_arctic_NDJF_overlay.png}
\end{subfigure}\\[0.5em]
% --- bottom row ---
\begin{subfigure}[b]{0.48\textwidth}
  \caption{}
  \includegraphics[width=\textwidth]{../figures/lwp_arctic_NDJF_overlay.png}
\end{subfigure}
\hfill
\begin{subfigure}[b]{0.48\textwidth}
  \caption{}
  \includegraphics[width=\textwidth]{../figures/twp_arctic_NDJF_overlay.png}
\end{subfigure}
\caption{Arctic-mean (60--90$^\circ$N) NDJF response as a function of INP
concentration for all MCT-MIP Level~0 models and configurations.
(a) Net cloud radiative effect (CRE, W\,m$^{-2}$); the control value is
shown at the left margin.
(b) Ice water path (IWP, g\,m$^{-2}$); E3SM's anomalously high and
monotonically increasing IWP explains its warming CRE response.
(c) Liquid water path (LWP, g\,m$^{-2}$).
(d) Total water path (TWP, g\,m$^{-2}$).
Error bars denote the inter-annual standard error of the mean.}
\label{fig:arctic_overlay}
\end{figure}

\paragraph{Level 1 protocol.}
As in Level~0, MCT-L1 simulations are atmosphere-only with prescribed
present-day SSTs and sea-ice fractions, with the INP perturbation implemented
as a transparent addition to background dust (modifying only freezing rates,
without a new tracer).
The perturbation range is revised to $10^2$, $10^4$, and
$10^6$\,L$^{-1}$ to better characterize the onset of the MCT response at low
concentrations.
In addition, participating models are encouraged to run a complementary batch
in which the background INP concentration from the control is multiplied by
$10^{-2}$, $10^{-4}$, and $10^{-6}$, probing the sensitivity to a lower INP
mean-state concentration.
It is expected that such experiments can help clarify how the INP mean-state
of models affects their MCT response.
\begin{table}[h]
\centering
\footnotesize
\begin{tabular}{llll}
\hline
\textbf{Group} & \textbf{Simulation} & \textbf{INP perturbation} & \textbf{Type} \\
\hline
Control        & CTRL                & none                                      & —              \\
\hline
\multirow{3}{*}{MCT-L1-Background}
               & INP\_low\_back      & CTRL $\times\,10^{-6}$\,L$^{-1}$        & multiplicative \\
               & INP\_mid\_back      & CTRL $\times\,10^{-4}$\,L$^{-1}$        & multiplicative \\
               & INP\_high\_back     & CTRL $\times\,10^{-2}$\,L$^{-1}$        & multiplicative \\
\hline
\multirow{3}{*}{MCT-L1-Seeding}
               & INP\_low\_seed      & CTRL $+\,10^{2}$\,L$^{-1}$              & additive       \\
               & INP\_mid\_seed      & CTRL $+\,10^{4}$\,L$^{-1}$              & additive       \\
               & INP\_high\_seed     & CTRL $+\,10^{6}$\,L$^{-1}$              & additive       \\
\hline
\end{tabular}
\caption{MCT-L1 simulation matrix (7 simulations total). MCT-L1-Seeding
perturbations add a spatially uniform INP concentration to the background dust;
MCT-L1-Background perturbations scale the control background by the given factor.}
\label{tab:mct_l1_sims}
\end{table}

The full set of 7 simulations (1 CTRL + 3 MCT-L1-Seeding +
3 MCT-L1-Background) is designed to map the complete range of
mixed-phase cloud responses to INP forcing and to reveal base-state biases of
the kind identified in Level~0.

\paragraph{Required output.}
Correctly attributing the MCT response requires cloud fields at higher
temporal resolution than standard monthly means; in particular, monthly-mean
bulk phase ratios are insufficient to reproduce satellite-based cloud-top
phase climatologies stratified by temperature.
In addition to the standard 2D monthly-mean fields from Level~0, the Level~1
protocol therefore requires:
\begin{itemize}
  \item \textbf{Daily sampled instantaneous values} of 2D fields (TOA, ATM,
        and BOA clear-sky and all-sky radiative fluxes) and 3D cloud fields
        (cloud water content and particle number for liquid and ice; snow and
        rain if available);
  \item \textbf{COSP forward-simulator} output (CALIPSO/ISCCP diagnostics)
        for direct comparison with satellite cloud-top-phase retrievals;
  \item \textbf{Glaciation process rates} (daily sampled instantaneous values):
        at least heterogeneous droplet freezing rate, ice depositional growth
        rate (WBF), and ice/snow riming rate.
        Because daily 3D process rates are very expensive to store, we encourage
        the use of online 2D column-integrated, mixed-phase-integrated, and
        cloud-top alternatives, which will otherwise be calculated offline.
\end{itemize}

\paragraph{Nudging.}
Wind-only nudging to ERA5 reanalysis is strongly encouraged to suppress
internal variability and enable meaningful multi-model comparison; free-running
simulations with a minimum of 10 years are an acceptable alternative.
